\section{Production model details}
\label{app:production}

\subsection{Pipeline, capacities and data}
\label{app:prod_setup}

The data are $336 \times 336$ pixel tissue-microarray cores collected by
quantum-cascade-laser infrared spectroscopy over the $1000$--$1800$~cm$^{-1}$
fingerprint region, with per-pixel four-class tissue labels. Each pixel carries
an $S = 942$ dimensional vector, obtained by stacking baseline-corrected
absorbance with its first two derivatives. The compositional model is the Block
Vision Transformer of \citet{oleary2026spatial}: a linear per-pixel
$f_\thetap : \R^{942} \to \R^{K}$, then $16 \times 16$ patch tokenization, a
12-layer, 12-head vision transformer of embedding width $M$, a
transposed-convolution decoder and a convolutional segmentation head. Parameter
counts in Table~\ref{tab:app_capacity} are read from the instantiated model.

\begin{table}[htbp]
\centering
\caption{Module capacities of the production pipeline, counted from the
instantiated model. $K$ is the spectral bottleneck width. The last row is the
single-channel field-standard input convention, for comparison.}
\label{tab:app_capacity}
\begin{tabular}{lrrr}
\toprule
Configuration & $\Cf$ & $\Cg$ & $\Cg/\Cf$ \\
\midrule
$K = 64$, three-channel input  & $61{,}108$  & $18{,}271{,}540$ & $299$ \\
$K = 128$, three-channel input & $121{,}588$ & $21{,}472{,}564$ & $177$ \\
$K = 64$, single-channel input & $20{,}916$  & $18{,}271{,}540$ & $874$ \\
\bottomrule
\end{tabular}
\end{table}

\paragraph{Channel convention.} The published architecture operates on the
second-derivative spectrum alone, the field-standard preprocessing in infrared
histopathology. The models analysed here instead stack raw absorbance with its
first two derivatives, following a companion tokenization study (anonymous,
under review). The choice affects only $\Cf$, and the three-channel convention
is the conservative one for the present purpose: under the field standard the
capacity asymmetry is roughly three times larger. Within the spatial module the
dominant blocks are the transformer ($5.34 \times 10^6$ parameters) and the
transposed convolution ($9.44 \times 10^6$), both $\Theta(M^2)$ in the
embedding width.

\paragraph{The head hypothesis is structurally unmet.} The segmentation head is
the chain $\mathrm{Conv}_{3\times3}(M{+}K \to 96) \to \mathrm{BN} \to
\mathrm{ReLU} \to \mathrm{Conv}_{3\times3}(96 \to 48) \to \mathrm{BN} \to
\mathrm{ReLU} \to \mathrm{Conv}_{3\times3}(48 \to 4)$, so the final affine
classifier sees $48 \times 3 \times 3 = 432$ input features at every width we
build ($M \in \{48, 192, 384\}$, $K \in \{64, 128\}$): the $\Theta(M)$ feature
map is separated from the logits by two nonlinearities. The
curvature-disparity theorem requires a dense classifier over penultimate
features of dimension $\Theta(M)$ with no downstream nonlinearity, and that
hypothesis is not satisfied here. Failing a hypothesis removes a guarantee; it
does not reverse the conclusion.

\subsection{Geometry at initialization}
\label{app:prod_geometry}

Gauss--Newton blocks of the instantiated pipeline were measured on real
training batches at random initialization, by Lanczos with full
reorthogonalization on matrix-free GGN block operators, validated to $10^{-12}$
against dense computation, sweeping $M \in \{48, 96, 192, 384\}$ with three
seeds and four batches per width. Table~\ref{tab:app_geom} collects the
results.

\begin{table}[htbp]
\centering
\small
\setlength{\tabcolsep}{4pt}
\caption{Production geometry at initialization. Ratios are mean $\pm$ sd over
three seeds and four batches per width. Slopes are log--log fits over the width
sweep. Effective ranks are out of $S = 942$. Directional entries are measured
at $M = 192$ over the nine measurements whose batch contains both classes
(three seeds $\times$ three batches).}
\label{tab:app_geom}
\begin{tabular}{lcc}
\toprule
Quantity & Breast & Prostate \\
\midrule
$\Dcurv$ at $M = 48$ & $0.33 \pm 0.12$ & --- \\
$\Dcurv$ at $M = 192$ & $0.46 \pm 0.21$ & --- \\
$\Dcurv$ at $M = 384$ & $0.70 \pm 0.19$ & $1.09 \pm 0.52$ \\
log--log slope of $\lambda_{\max}(G_{\thetap\thetap})$ vs $M$ & $-0.003$ & $-0.012$ \\
log--log slope of $\lambda_{\max}(G_{\phip\phip})$ vs $M$ & $0.38$ ($115 \to 258$) & $0.45$ \\
Effective rank of $\Sigma_X$, raw inputs & $1.07$ & $1.02$ \\
Effective rank of $\Sigma_X$, post-normalization & $1.21$ & --- \\
$|\cos|$ between $v_1$ and the mean spectrum (post-norm.\,/\,raw) & $>0.995$\,/\,$>0.9999$ & --- \\
$\thetap$-block top eigenvector's norm fraction in the $v_1$-induced subspace & $0.954 \pm 0.017$ & --- \\
Trace fraction in the top $40$ of $61{,}108$ spectral directions & $0.83 \pm 0.12$ & --- \\
Curvature ratio, $v_1$ vs full class contrast ($|\cos(d,v_1)| = 0.45$) & $3.7\times$ ($2.6$--$4.8$) & --- \\
Curvature ratio, $v_1$ vs $v_1$-orthogonalized contrast & $12$--$28\times$ & --- \\
Squared-norm fraction of the contrast retained after orthogonalization & $76$--$84\%$ & --- \\
Rank at which $\lambda_{\max}(G_{\phip\phip})$ overtakes the $\thetap$ spectrum & $2$--$5$ & --- \\
\bottomrule
\end{tabular}
\end{table}

Two controls bracket the denominator. Removing the spectral normalization
deflates $\lambda_{\max}(G_{\thetap\thetap})$ by $2.4\times$ ($471 \to 194$ at
$M = 192$) and pushes the ratio across parity ($1.18$ at $M = 192$, $1.31$ at
$M = 384$). Input \emph{centering} changes nothing when the architecture
contains internal batch normalization: $\Dcurv$ is identical to three decimals
with and without centering, because a linear projection followed by
normalization absorbs constant input shifts exactly. The near-rank-one input
statistics make the cap side of the curvature-disparity bound uninformative at
buildable widths independently of the head-hypothesis failure: the data
concentrate almost all their second-moment mass in one direction, and that
direction is the mean spectrum. The directional measurements establish an
anisotropy, not its consequence for learning; whether it starves learning is a
hypothesis, not a measurement.

\subsection{What carries the width trend}
\label{app:prod_width}

The measurements in this subsection are made on the production model at random
initialization, at three widths $M \in \{48, 192, 384\}$ and three seeds, on
one core (the densest of the first $24$ breast fold-0 training cores, with
$29{,}664$ labelled pixels of $112{,}896$), with cuDNN TF32 disabled. Four
normalization modes are compared: \texttt{train} (both head batch
normalizations using batch statistics), \texttt{eval\_bn1} and
\texttt{eval\_bn2} (only the first or only the second head normalization at its
initial running statistics) and \texttt{eval\_head} (both). Artifacts:
\texttt{results/exp1\_8c\_interior\_witness.csv} and
\texttt{results/exp1\_8d\_REPORT.md}.

\paragraph{The incoming features are flat in width.} The only downstream layer
whose input dimension grows with $M$ is the first head convolution. The top
eigenvalue of the patch-unfolded second-moment matrix of the features entering
it is unchanged to four significant figures across the three widths, at
$356.3$, $403.7$ and $332.4$ for seeds $0$, $1$ and $2$ respectively; the
leading direction sits inside the width-independent $K = 64$ skip channels.
Table~\ref{tab:app_gram} also records a discrepancy that matters for any
normalization argument: the same Gram computed over the full batch-normalization
group (all positions, ignored pixels included) rather than over labelled pixels
is substantially smaller, and the cosine between the two centered top
eigenvectors is far from one.

\begin{table}[htbp]
\centering
\small
\caption{Incoming $3\times3$-patch Gram of the first head convolution, top
eigenvalue, from \texttt{results/exp1\_8d\_REPORT.md}. Values are identical
across $M \in \{48, 192, 384\}$ and across all four normalization modes, so a
single column suffices. \emph{labelled} uses the labelled-pixel mask;
\emph{full group} uses the batch-normalization group. The last column is the
cosine between the two centered top eigenvectors.}
\label{tab:app_gram}
\begin{tabular}{rrrrr}
\toprule
seed & labelled & labelled, centered & full group & $\cos$(labelled, full) \\
\midrule
$0$ & $356.3$ & $68.19$ & $263.3$ & $0.2793$ \\
$1$ & $403.7$ & $54.60$ & $306.8$ & $0.5502$ \\
$2$ & $332.4$ & $67.14$ & $253.0$ & $0.2176$ \\
\bottomrule
\end{tabular}
\end{table}

\paragraph{Normalization quantities.} Table~\ref{tab:app_bn} gives the
per-channel batch statistics of the two head normalizations, measured over the
group they actually normalize, together with the activation energy entering and
leaving the first one. Under fan-in initialization the pre-normalization
variance of the first head normalization falls roughly as $1/(M+K)$, and its
inverse-standard-deviation gain grows correspondingly. No channel's variance is
within a factor of ten of the stabilizer $\varepsilon = 10^{-5}$: the
fraction of channels with variance below $10\varepsilon$ is $0$ in every row
of both normalizations at every width, seed and mode. (This recorded check
does not establish the stronger separation by three orders of magnitude
claimed in the earlier draft.)

\begin{table}[htbp]
\centering
\small
\caption{Head normalization quantities at initialization, \texttt{train} mode,
from \texttt{results/exp1\_8d\_REPORT.md} Table 1. Per-seed values are
per-channel medians (variance, gain) or labelled-pixel energies; the
\emph{mean} column is the mean of the three seed values, computed by us.
$\mathrm{gain}_i = \gamma_i/\sqrt{\mathrm{var}_i + \varepsilon}$. The last row
is the product of the mean variance with $M + K$, $K = 64$.}
\label{tab:app_bn}
\begin{tabular}{lrrrrr}
\toprule
Quantity & seed $0$ & seed $1$ & seed $2$ & mean & $M$ \\
\midrule
First head BN, median variance & $0.1470$ & $0.1467$ & $0.1587$ & $0.1508$ & $48$ \\
                               & $0.0678$ & $0.0589$ & $0.0656$ & $0.0641$ & $192$ \\
                               & $0.0346$ & $0.0310$ & $0.0398$ & $0.0351$ & $384$ \\
First head BN, median gain     & $2.609$ & $2.611$ & $2.511$ & $2.577$ & $48$ \\
                               & $3.842$ & $4.120$ & $3.905$ & $3.956$ & $192$ \\
                               & $5.375$ & $5.682$ & $5.010$ & $5.356$ & $384$ \\
Second head BN, median gain    & $2.892$ & $3.122$ & $3.039$ & $3.018$ & $48$ \\
                               & $3.110$ & $3.177$ & $3.195$ & $3.161$ & $192$ \\
                               & $2.966$ & $3.060$ & $3.121$ & $3.049$ & $384$ \\
Pre-BN energy, labelled pixels & $25.64$ & $26.28$ & $23.91$ & $25.28$ & $48$ \\
                               & $11.07$ & $10.67$ & $11.19$ & $10.98$ & $192$ \\
                               & $5.765$ & $5.543$ & $6.015$ & $5.774$ & $384$ \\
Post-BN energy, labelled pixels & $132.1$ & $130.8$ & $133.2$ & $132.0$ & $48$ \\
                                & $130.4$ & $131.5$ & $132.5$ & $131.5$ & $192$ \\
                                & $130.7$ & $128.5$ & $130.8$ & $130.0$ & $384$ \\
\midrule
Mean variance $\times\,(M{+}K)$ & --- & --- & --- & $16.9\,/\,16.4\,/\,15.7$ & $48\,/\,192\,/\,384$ \\
\bottomrule
\end{tabular}
\end{table}

\paragraph{Witness Jacobian-vector products.} A unit perturbation of the first
head convolution is propagated by forward-mode differentiation through the
actual head in the actual normalization mode, so that the train-mode derivative
includes the dependence on the batch statistics. The final scalar
$e_{\mathrm{final}} = N^{-1}\sum_n d_n^\top(\mathrm{diag}(p_n) - p_np_n^\top)d_n$
is a lower-bound witness for $\lambda_{\max}(G_{\phip\phip})$ and is verified
against the same matvec the eigensolver uses; all $216$ rows pass at relative
discrepancy below $10^{-3}$, with a maximum of $1.269 \times 10^{-5}$ over the
$198$ rows with nonzero curvature. Table~\ref{tab:app_jvp} gives the stage
energies for two directions in two modes.

\begin{table}[htbp]
\centering
\small
\caption{Witness JVP stage energies, mean over three seeds, from
\texttt{results/exp1\_8d\_REPORT.md} Table 2a. \texttt{ggn\_top} is the full
top GGN eigenvector of the convolution block; \texttt{s9cen} is the centered
leading direction of the incoming patch Gram. $e_{\mathrm{pre}}$ is the
perturbation energy before the first head normalization,
$e_{\mathrm{logit}}$ is measured on labelled pixels.}
\label{tab:app_jvp}
\begin{tabular}{llrrrrrr}
\toprule
Mode & Direction & $M$ & $e_{\mathrm{pre}}$ & $e_{\mathrm{bn1}}$ & $e_{\mathrm{bn2}}$ & $e_{\mathrm{logit}}$ & $e_{\mathrm{final}}$ \\
\midrule
\texttt{train} & \texttt{ggn\_top} & $48$  & $260.6$ & $2875$ & $4453$ & $261.7$ & $60.73$ \\
\texttt{train} & \texttt{ggn\_top} & $192$ & $259.4$ & $5905$ & $7553$ & $505.2$ & $118.5$ \\
\texttt{train} & \texttt{ggn\_top} & $384$ & $256.8$ & $9168$ & $13810$ & $1018$ & $227.6$ \\
\texttt{train} & \texttt{s9cen}    & $48$  & $76.53$ & $1538$ & $1224$ & $29.23$ & $5.761$ \\
\texttt{train} & \texttt{s9cen}    & $192$ & $76.53$ & $2305$ & $1914$ & $50.83$ & $9.862$ \\
\texttt{train} & \texttt{s9cen}    & $384$ & $76.53$ & $3997$ & $3906$ & $124.3$ & $21.95$ \\
\midrule
\texttt{eval\_head} & \texttt{ggn\_top} & $48$  & $261.9$ & $261.9$ & $42.42$ & $9.607$ & $2.366$ \\
\texttt{eval\_head} & \texttt{ggn\_top} & $192$ & $262.8$ & $262.8$ & $42.18$ & $10.35$ & $2.487$ \\
\texttt{eval\_head} & \texttt{ggn\_top} & $384$ & $262.5$ & $262.5$ & $46.53$ & $10.45$ & $2.372$ \\
\texttt{eval\_head} & \texttt{s9cen}    & $48$  & $76.53$ & $76.53$ & $9.103$ & $0.4667$ & $0.1049$ \\
\texttt{eval\_head} & \texttt{s9cen}    & $192$ & $76.53$ & $76.53$ & $8.460$ & $0.4613$ & $0.1013$ \\
\texttt{eval\_head} & \texttt{s9cen}    & $384$ & $76.53$ & $76.53$ & $9.762$ & $0.4450$ & $0.1011$ \\
\bottomrule
\end{tabular}
\end{table}

A spatially constant bias perturbation $\Delta b = e_c$ is annihilated exactly
by the immediately following train-mode normalization: in \texttt{train} and
\texttt{eval\_bn2} modes, at every width and seed, all stage energies after the
centering step are exactly zero and the quadratic form agrees with the matvec
to at most $1.14\times10^{-17}$ in absolute value. In \texttt{eval\_head} the
same perturbation gives $e_{\mathrm{final}}$ between $7.8\times10^{-4}$ and
$2.0\times10^{-3}$, and in \texttt{eval\_bn1} between $0.0154$ and $0.0692$.

\paragraph{Slopes and attribution.} Table~\ref{tab:app_slopes} gives the
log--log slopes against $M$ per seed and pooled, together with the
mode-to-mode ratios of levels.

\begin{table}[htbp]
\centering
\small
\caption{Log--log slopes against $M$ over the three widths, per seed and
pooled, and mode-to-mode ratios of the levels (each mode's value divided by its
train-mode value at the same width and seed, then averaged over seeds and over
the three widths), from
\texttt{results/exp1\_8d\_REPORT.md}. $\lambda_{WW}$ is the top curvature of
the first head convolution's weight block; $\lambda_\phip$ is the whole spatial
block. Three-width slopes are descriptive, not asymptotic exponents.}
\label{tab:app_slopes}
\begin{tabular}{lrrrrrrr}
\toprule
& \multicolumn{4}{c}{slope of $\lambda_{WW}$} & \multicolumn{1}{c}{$\lambda_\phip$} & \multicolumn{2}{c}{level ratio vs \texttt{train}} \\
\cmidrule(lr){2-5}\cmidrule(lr){6-6}\cmidrule(lr){7-8}
Mode & seed $0$ & seed $1$ & seed $2$ & pooled & pooled slope & $\lambda_{WW}$ & $\lambda_\phip$ \\
\midrule
\texttt{train}     & $0.509$ & $0.785$ & $0.350$ & $+0.613$ & $+0.398$ & $1$ & $1$ \\
\texttt{eval\_bn1} & $0.429$ & $0.634$ & $0.542$ & $+0.563$ & $+0.281$ & $0.903$ & $1.03$ \\
\texttt{eval\_bn2} & $0.544$ & $0.632$ & $0.508$ & $+0.571$ & $+0.253$ & $0.156$ & $0.227$ \\
\texttt{eval\_head} & $-0.112$ & $0.181$ & $-0.084$ & $+0.006$ & $-0.331$ & $0.026$ & $0.037$ \\
\bottomrule
\end{tabular}
\end{table}

Three readings, and one non-reading. The positive width trend in the spatial
block's curvature at initialization is sensitive to head batch normalization:
holding both head normalizations at their initial running statistics removes it
($+0.613 \to +0.006$ for the convolution block, $+0.398 \to -0.331$ for the
whole spatial block) while leaving the incoming features unchanged. Either
normalization alone preserves the trend ($+0.563$ and $+0.571$), which is
evidence against attributing it uniquely to the first one, even though most of
the drop in the \emph{level} attaches to the second. The mechanism this
supports is a finite-range normalization gain acting on width-diluted
activations, over the range of widths measured, and not the width-linear
feature-Gram mechanism of the shallow theory. The non-reading: the
\texttt{eval\_head} intervention changes the normalization derivative, the
forward logits, the downstream gates and the cross-entropy metric at the same
time, so a difference between modes establishes sensitivity to the
intervention, not a unique cause; and the observation that the variance never
approaches the stabilizer is a heuristic about the regime, not a guarantee.

\paragraph{Parameterization control.} Table~\ref{tab:app_scale} scales the
weights and bias of the first head convolution by $c$ and, in the exact
invariance, the following stabilizer by $c^2$. The logits are unchanged and the
block's top curvature scales by exactly $c^{-2}$: the raw Euclidean curvature
of a layer feeding a normalization is a parameterization quantity, not an
architectural one \citep{vanlaarhoven2017l2}.

\begin{table}[htbp]
\centering
\small
\caption{Function-preserving scale control at seed $0$, from
\texttt{results/exp1\_8d\_REPORT.md} Table 3. \emph{ratio} is
$c^2\lambda_{WW}(c)/\lambda_{WW}(1)$, which is $1$ under exact invariance.
\emph{logit dev.} is $\max|\text{logits}(c) - \text{logits}(1)|$, reported both
with the stabilizer co-scaled by $c^2$ and with it held fixed (the
implementation's actual behaviour); $\lambda_{WW}$, \emph{ratio} and
$\norm{W}_F^2$ are taken from the fixed-stabilizer variant, and the co-scaled
variant differs from it by at most one unit in the fourth significant figure.}
\label{tab:app_scale}
\begin{tabular}{rrrrrrr}
\toprule
$M$ & $c$ & $\lambda_{WW}$ & ratio & logit dev.\ (co-scaled) & logit dev.\ (fixed $\varepsilon$) & $\norm{W}_F^2$ \\
\midrule
$48$  & $1/4$ & $719.4$ & $1.000$ & $0$ & $9.40\cdot10^{-4}$ & $2.003$ \\
$48$  & $1/2$ & $179.8$ & $1.000$ & $0$ & $1.89\cdot10^{-4}$ & $8.014$ \\
$48$  & $1$   & $44.96$ & $1.000$ & $0$ & $0$ & $32.06$ \\
$48$  & $2$   & $11.24$ & $1.000$ & $0$ & $4.63\cdot10^{-5}$ & $128.2$ \\
$48$  & $4$   & $2.810$ & $1.000$ & $0$ & $5.84\cdot10^{-5}$ & $512.9$ \\
$192$ & $1/4$ & $1935$  & $0.9998$ & $0$ & $3.00\cdot10^{-3}$ & $2.002$ \\
$192$ & $4$   & $7.560$ & $1.000$ & $0$ & $1.89\cdot10^{-4}$ & $512.6$ \\
$384$ & $1/4$ & $1927$  & $0.9991$ & $0$ & $4.69\cdot10^{-3}$ & $1.999$ \\
$384$ & $4$   & $7.535$ & $1.000$ & $0$ & $2.95\cdot10^{-4}$ & $511.8$ \\
\bottomrule
\end{tabular}
\end{table}

\paragraph{Eigensolver accuracy.} Every row of the width sweep terminates on
the combined criterion (relative change of the Rayleigh quotient below
tolerance \emph{and} residual below $10^{-3}$), never on the iteration or
wall-clock budget. Over all $36$ (mode, width, seed) rows the relative residual
$\norm{Gv - \lambda v}/\lambda$ is at most $3.81\times10^{-4}$ for the
convolution weight block, at most $3.81\times10^{-4}$ for the weight--bias
block and at most $9.97\times10^{-4}$ for the whole spatial block, and the
Rayleigh quotient of the returned iterate agrees with the reported eigenvalue
to the printed precision in every row.

\subsection{Controlled retraining: original protocol}
\label{app:prod_e3c}

The pipeline was retrained on breast fold $0$ at reduced depth (six transformer
layers) with full instrumentation. Two protocols were run and both are
reported: the \emph{original} protocol ($39$ runs), which is the pre-registered
one and whose joint-versus-frozen contrast an audit of the training code later
showed to be confounded, and the \emph{matched} protocol ($21$ runs), which
removes the confounds, saves the selected checkpoint and re-runs the
pre-registered comparisons. Sixty logged runs in total. Per-step residuals,
block gradient norms, per-epoch parameter displacements and input-Jacobian
norms are recorded in both; the matched protocol additionally logs the
per-step clipping coefficient and the per-block applied update norms.

The original protocol comprises arms $\{$joint-linear, frozen-random,
frozen-PCA, joint-MLP$\}$ $\times$ widths $M \in \{48, 192\}$ $\times$ three
seeds under AdamW for $60$ epochs, plus a $\{$joint-linear, frozen-random$\}$
$\times$ three-seed pair under momentum SGD at both widths and a frozen-PCA SGD
triple at $M = 48$ as a positive control. The spectral parameters sit in a
zero-weight-decay group in every arm, and frozen arms log counterfactual
spectral gradients. Table~\ref{tab:app_e3c_orig} gives the AdamW arms.

\begin{table}[htbp]
\centering
\small
\caption{Controlled study, original protocol, AdamW arms: validation macro-F1
(mean $\pm$ sd over three seeds) under two checkpoint policies, and relative
spectral displacement at end of training. \emph{Final-5} is the pre-registered
metric. \emph{Best-val} is the maximum over the $60$ per-epoch evaluations on
the same $28$ validation cores; no checkpoint was saved at that maximum and
this study has no test set, so that column is exploratory and
selection-optimistic. The two policies order the arms oppositely. These arms
also differ in clip scope and in normalization-affine treatment, so this
contrast is confounded.}
\label{tab:app_e3c_orig}
\begin{tabular}{lccccr}
\toprule
& \multicolumn{2}{c}{Final-5 (pre-registered)} & \multicolumn{2}{c}{Best-val (exploratory)} & \\
\cmidrule(lr){2-3}\cmidrule(lr){4-5}
Arm & $M{=}48$ & $M{=}192$ & $M{=}48$ & $M{=}192$ & $\norm{\Delta\thetap}/\norm{\thetap_0}$ \\
\midrule
Joint (linear) & $0.538 \pm 0.082$ & $0.646 \pm 0.050$ & $0.873 \pm 0.030$ & $0.858 \pm 0.056$ & $0.54$--$0.64$ \\
Frozen random  & $0.648 \pm 0.063$ & $0.731 \pm 0.023$ & $0.840 \pm 0.025$ & $0.817 \pm 0.017$ & $0$ \\
Frozen PCA     & $0.731 \pm 0.003$ & $0.726 \pm 0.012$ & $0.769 \pm 0.021$ & $0.862 \pm 0.068$ & $0$ \\
Joint (MLP)    & $0.566 \pm 0.070$ & $0.575 \pm 0.184$ & $0.900 \pm 0.000$ & $0.897 \pm 0.001$ & $0.85$--$0.96$ \\
\bottomrule
\end{tabular}
\end{table}

On the pre-registered final-window metric the joint-linear minus
frozen-random gaps are $-0.110$ at $M = 48$ ($90\%$ paired-$t$ interval
$[-0.333, +0.114]$) and $-0.085$ at $M = 192$ ($[-0.185, +0.015]$), neither
resolvable at $n = 3$; against frozen-PCA they are $-0.193$
($[-0.334, -0.052]$) and $-0.081$ ($[-0.175, +0.014]$). Taking each run's
maximum over its $60$ evaluations reverses the ordering: joint minus
frozen-random becomes $+0.033$ and $+0.041$, joint minus frozen-PCA $+0.104$
and $-0.004$. The switch is an arithmetic identity --- the change in every gap
equals the difference in peak-to-final degradation between the two arms ---
and averaged over seeds that degradation is $0.335$ ($M = 48$) and $0.212$
($M = 192$) for joint-linear against $0.192$ and $0.086$ for frozen-random.

\paragraph{Three asymmetries.} First, gradient-norm clipping at $1.0$ was
computed over the optimizer's parameters, that is over $\thetap + \phip$ in the
joint arms but over $\phip$ alone in the frozen arms, so at identical gradients
the frozen arm takes the larger $\phip$ step: the median extra clip factor paid
by the joint arm is $1.10$--$1.12$ under SGD at $M = 48$ and $3.0$--$3.6$ under
AdamW at $M = 192$. Second, the normalization following the spectral projection
contributes $128$ affine parameters that are trained in joint-linear, frozen in
frozen-random and structurally absent in frozen-PCA. Third, no checkpoint was
saved during training, so the best-validation column is a post-hoc curve
maximum that no saved model realizes, taken over $60$ evaluations of the same
$28$ validation cores that define the comparison.

\subsection{Momentum-SGD controls and normalization recalibration}
\label{app:prod_sgd}

The two optimizer families constitute an optimizer-family control; the SGD arm
is not an instantiation of the two-timescale theorem's gradient flow. It uses
momentum $0.9$; weight decay $0.01$ on $\phip$ and $0$ on $\thetap$;
mini-batches of four cores with reshuffling; and global gradient-norm clipping
at $1.0$, which binds on $73$--$89\%$ of steps. Logged gradient norms are
recorded before clipping, so the plotted theory quantities are not the
quantities the trajectory applied. Under momentum SGD the paired joint minus
frozen-random gap is $+0.002$ at $M = 48$ and $+0.011$ at $M = 192$ on
best-validation, with $90\%$ paired-$t$ intervals $[-0.008, +0.011]$ and
$[-0.007, +0.030]$; on the final-window metric the same pairs give $+0.028$ and
$-0.057$, with intervals $[-0.173, +0.228]$ and $[-0.124, +0.010]$ --- gaps
that differ in sign between the widths, with intervals an order of magnitude
wider than the effect they bracket. Both arms lose $0.12$--$0.22$ macro-F1
between their peak and their final window. Descriptively, the spectral share of
the summed squared gradient norms drops from about $0.94$ (AdamW joint) to
about $0.26$ (SGD joint); momentum, clipping, weight decay and adaptivity all
differ between the two families, so no single factor is isolated.

\paragraph{Recalibration of normalization statistics.} On disposable copies of
the final checkpoints, with all weights fixed, the normalization running
statistics were re-estimated on training data only ($29$ batches) and the model
re-evaluated. For this intervention and these checkpoints, the recovered
fraction of the joint arms' peak-to-final drop is $0.17$ at $M = 48$ and
$0.07$ at $M = 192$: recalibration recovers a nonzero part of the degradation
but not most of it. What it does undo is the instability of the endpoint
number: it repairs the two collapsed joint checkpoints ($+0.24$ and $+0.15$
macro-F1) and shrinks the seed standard deviation by roughly $3.7\times$, while
leaving the frozen arms essentially unchanged ($|\Delta| \le 0.02$, and
$\le 0.006$ for frozen-PCA, whose spectral stage carries no normalization).

\subsection{Matched protocol}
\label{app:prod_matched}

The matched protocol equalizes the three asymmetries and re-runs the
pre-registered comparisons at $M = 48$: the reduction-stage normalization
affine parameters are placed in $\phip$ and trained wherever that
normalization exists; the clip norm at $1.0$ is computed over $\phip$ only
under the same clipping rule and scope in every arm; and the best-validation
checkpoint is saved. AdamW at learning rate $10^{-4}$, $60$ epochs, three
seeds, flip and $90^\circ$-rotation augmentation on in all arms. The logged
per-step clip coefficient averages $0.33$--$0.57$ across arms --- the
coefficients and gradients differ by arm, which is an outcome of the matched
rule --- and the logged applied update norms confirm that the spectral share of
the applied update tracks the learning-rate multiplier ($0.02$, $0.17$ and
$1.36$ times the $\phip$ update norm for $\times0.1$, $\times1$ and
$\times10$). Table~\ref{tab:app_e3c_matched} gives all seven arms.

\begin{table}[htbp]
\centering
\small
\caption{Controlled study, matched protocol ($M = 48$, AdamW, three seeds,
mean $\pm$ sd). Identical normalization-affine treatment and the same
$\phip$-only clipping rule and scope in every arm; best-validation checkpoints
saved. \emph{Final-5} is the pre-registered metric; \emph{best-val} is the
maximum over $60$ per-epoch evaluations on the $28$ fold-0 validation cores and
is exploratory and selection-optimistic, since this study has no test set. The
drift column is the relative spectral-parameter displacement at end of
training; its standard deviations were computed by us from
\texttt{results/e3c\_analysis\_runs.csv}. Rows $1$--$2$ and rows $3$--$4$ are
the two matched pairs; see the text.}
\label{tab:app_e3c_matched}
\begin{tabular}{lccc}
\toprule
Arm & Final-5 (pre-registered) & Best-val (exploratory) & $\norm{\Delta\thetap}/\norm{\thetap_0}$ \\
\midrule
Frozen random                   & $0.684 \pm 0.037$ & $0.814 \pm 0.040$ & $0$ \\
Joint (linear)                  & $0.694 \pm 0.056$ & $0.837 \pm 0.029$ & $0.566 \pm 0.026$ \\
Frozen pretrained               & $0.690 \pm 0.033$ & $0.861 \pm 0.022$ & $0$ \\
Fine-tuned pretrained           & $0.710 \pm 0.052$ & $0.857 \pm 0.009$ & $0.020 \pm 0.005$ \\
Joint, spectral lr $\times 0.1$ & $0.626 \pm 0.055$ & $0.841 \pm 0.007$ & $0.083 \pm 0.006$ \\
Joint, spectral lr $\times 10$  & $0.561 \pm 0.114$ & $0.886 \pm 0.037$ & $4.910 \pm 0.402$ \\
Joint, cosine schedule          & $0.660 \pm 0.020$ & $0.825 \pm 0.033$ & $0.359 \pm 0.013$ \\
\bottomrule
\end{tabular}
\end{table}

\paragraph{Two matched pairs, not a single-factor design.} The four core arms
do not all differ in one factor. Joint-linear and frozen-random use a linear
reduction stage with normalization; frozen-pretrained and fine-tuned-pretrained
use a pretrained multilayer spectral encoder without reduction-stage
normalization. The clean freezing comparisons are therefore the two matched
pairs, joint-linear against frozen-random and fine-tuned against
frozen-pretrained. Comparing across the pairs --- frozen-pretrained against
frozen-random, for instance --- changes architecture, initialization and
training history, and validation selection, not only informativeness.

\paragraph{Verdicts.} All comparisons in Table~\ref{tab:app_verdicts} are
paired by seed with $90\%$ $t$-intervals ($t = 2.920$, $n = 3$), conditional on
one validation split, and are descriptive seed summaries rather than a
population or multi-factor causal analysis.

\begin{table}[htbp]
\centering
\small
\setlength{\tabcolsep}{4pt}
\caption{Matched protocol: paired differences with $90\%$ paired-$t$ intervals
($n = 3$). The last two rows compare the two protocols and are reported as
descriptive seed summaries; saving a checkpoint changes what is available to
evaluate, not the final-window trajectory.}
\label{tab:app_verdicts}
\begin{tabular}{lcc}
\toprule
Contrast & Final-5 & Best-val \\
\midrule
Fine-tuned $-$ frozen-pretrained & $+0.021\ [-0.022, +0.063]$ & $-0.004\ [-0.048, +0.039]$ \\
Joint-linear $-$ frozen-random & $+0.011\ [-0.142, +0.163]$ & $+0.023\ [-0.066, +0.112]$ \\
Frozen-pretrained $-$ frozen-random & --- & $+0.047\ [-0.042, +0.136]$ \\
Spectral lr $\times10 - \times1$ & $-0.133\ [-0.419, +0.152]$ & $+0.048\ [-0.016, +0.112]$ \\
Spectral lr $\times0.1 - \times1$ & $-0.068\ [-0.173, +0.037]$ & $+0.004\ [-0.039, +0.046]$ \\
\midrule
Joint/frozen gap: matched $-$ original & $+0.120\ [-0.230, +0.470]$ & $-0.010\ [-0.105, +0.085]$ \\
\bottomrule
\end{tabular}
\end{table}

The pre-registered prediction that fine-tuning a pretrained spectral encoder
falls below leaving it frozen is not supported: neither difference is
resolvable at $n = 3$, and the fine-tuned arm moved $\thetap$ by only $2\%$ of
its initial norm. The logged applied update norms verify nonzero applied
encoder updates under AdamW in that arm; this is a statement about the applied
updates, not an attribution of the loss decrease. The pre-registered
harmful-drift ordering $\times0.1 \ge \times1 \ge \times10$ is not observed: on
best-validation the $\times10$ arm is above $\times1$ in $3/3$ seeds while its
final-window endpoint is the lowest and most variable of any arm at $0.561$.
Three multiplier means do not establish a systematic relation between
displacement and peak score, and the intervals are wide. The cosine schedule
improved neither mean endpoint ($0.825$ best-val and $0.660$ final-5 against
$0.837$ and $0.694$ for the constant learning rate); one schedule that did not
help does not rule out an excessive late learning rate.

\paragraph{Disclosures.} Width $48$ only; $n = 3$; a single validation fold of
$28$ fold-0 cores; no test set. The best-validation column is exploratory and
the final-window column is the pre-specified metric. The pretrained encoders
were selected by pixel validation macro-F1 on the \emph{same} fold-0 validation
cores the runner later scores, which is a selection advantage for both
pretrained arms, and were pretrained with weight decay $0.01$ on $\thetap$.
Pixel pretraining used the unaugmented centre-crop protocol while subsequent
training used random crop plus augmentation; this affects comparisons across
the two pairs, not the frozen-versus-fine-tuned pair that starts from the same
checkpoint. Augmentation is on in all arms. Finally, the summary artifact
\texttt{results/e3c\_analysis\_runs.csv} does not expose per-run configuration
hashes or data-order pairing, so before the pairs are described as completely
matched a compact provenance manifest should accompany it, recording
architecture, initial checkpoint hash, parameter-group membership, clipping and
weight-decay settings, normalization mode and state, augmentation and
data-order seed, scheduler and selection rule; identical seed labels alone do
not establish bit-identical mini-batch order across different model
constructions.

\subsection{What is not certified on this pipeline}
\label{app:prod_scope}

The spatial module's input-Jacobian operator norm grows over training, from
$3.2$ to $157$ on average for joint-linear (up to $414$) and from $1.7$ to
$1.2\times10^{4}$ for joint-MLP, which rejects an initialization-scale
stability approximation for the containment constant of the displacement
theorem, though not the existence of a finite-horizon constant over the
training window. The fitted per-step residual-envelope rate for joint-linear
grows from $\hat\mu = 6.0\times10^{-4}$ ($M = 48$) to $1.9\times10^{-3}$
($M = 192$), a factor $3.2$ for a $4\times$ width increase; the ratio is
invariant to the step-clock-to-flow-time conversion, the absolute values are
not, and two widths do not establish a trend.

